\subsection{Hausdorff continuity does not supply face parameterizations}

The topology on unparameterized outer faces cannot be replaced silently
by the Hausdorff topology.  The following example separates these two
notions even for semialgebraic polygonal disks.

\begin{example}[A collapsing hairpin]\label{ex:hausdorff-hairpin}
In the square $X=[-3,3]^2$, let $C_0=[-2,2]^2$ and, for
$0<\varepsilon\leq 1/4$, put
\[
 C_\varepsilon=[-2,2]^2
 \cup\bigl([1,2]\times[2,2+\varepsilon]\bigr)
 \cup\bigl([-1,2]\times[2+\varepsilon,2+2\varepsilon]\bigr).
\]
Every $C_\varepsilon$ is a PL disk containing the origin in its interior,
so it is a strict cone neighborhood with a collared outer face.
Both $C_\varepsilon$ and $\partial C_\varepsilon$ converge in Hausdorff
distance to $C_0$ and $\partial C_0$, respectively.  Nevertheless, no
neighborhood of $0$ in the parameter interval admits a jointly
continuous family of homeomorphisms
\[
 e_\varepsilon:S^1\xrightarrow{\sim}\partial C_\varepsilon.
\]
Over an arbitrary real closed field the same conclusion holds for
jointly definable continuous parameterizations.
\end{example}

\begin{proof}
The two added rectangles attach successively along boundary intervals.
Their union with the square is therefore a PL disk.  Equivalently, its
boundary is the simple polygon with the following cyclic list of
vertices:
\[
 \begin{gathered}
 (-2,2),\ (1,2),\ (1,2+\varepsilon),\ (-1,2+\varepsilon),\\
 (-1,2+2\varepsilon),\ (2,2+2\varepsilon),\ (2,-2),\ (-2,-2).
 \end{gathered}
\]
A PL disk with a specified interior point is a PL cone on its boundary
with that point as vertex.  Its polygonal boundary also has an exterior
collar inside $X$.  All changed points lie in the strip
$2\leq y\leq2+2\varepsilon$.  Vertical projection, together with the
unchanged sides, shows that both Hausdorff distances are at most
$2\varepsilon$.  For positive parameters, a piecewise linear increasing
change of the $y$ coordinate identifies the disks and their boundaries,
so the only issue is the limit at zero.

Suppose a joint parameterization $e$ exists on $[0,\delta]\times S^1$.
Set
\[
 P_\varepsilon=(0,2),\qquad
 Q_\varepsilon=(0,2+\varepsilon),\qquad P=(0,2).
\]
Both marked points belong to $\partial C_\varepsilon$ and converge to
$P$.  If $a=e_0^{-1}(P)$, compactness of the circle and injectivity of
$e_0$ imply
\[
 e_\varepsilon^{-1}(P_\varepsilon)\longrightarrow a,
 \qquad
 e_\varepsilon^{-1}(Q_\varepsilon)\longrightarrow a.
\]
Choose a small closed arc $I$ about $a$ whose image under $e_0$ lies in
the ball of radius $1/4$ about $P$.  Joint continuity implies that, for
small positive $\varepsilon$, the whole arc $e_\varepsilon(I)$ lies in
the ball of radius $1/2$ about $P$ and contains both marked points.
The subarc joining them is one of the two arcs of the polygonal
boundary.  The first passes through $(1,2)$, while the other passes
through $(-1,2+\varepsilon)$.  Each contains a point at distance at
least one from $P$, a contradiction.

For the definable version, the map
\[
 [0,\delta]\times S^1\longrightarrow
 \{(\varepsilon,z):z\in\partial C_\varepsilon\},
 \qquad(\varepsilon,x)\longmapsto(\varepsilon,e_\varepsilon(x))
\]
would be a continuous definable bijection from a definably compact set
to a definable Hausdorff set.  Its inverse is continuous.  This gives
the same convergence of the two preimages, and definable compactness
gives the uniform control on $e_\varepsilon(I)$.  The polygonal arc
contradiction uses only the ordered-field inequalities above.
\end{proof}

Thus a result about outer-face families with local joint
parameterizations does not by itself answer a question about
Hausdorff-continuous families.  In particular, this example makes no
claim about whether a prescribed six-torus loop can be filled in the
Hausdorff-only space.
